\chapter{寄存器、计数器与状态机}

一个触发器保存 1 bit。把多个触发器并排，并让它们共享同一个时钟边界，就得到多 bit 寄存器。再把寄存器输出送进组合逻辑，把组合逻辑结果送回寄存器输入，就能让硬件按规则一步一步改变状态。

\section{一、寄存器保存一组 bit}

寄存器可以理解为一组共享时钟和写使能的触发器。若写使能有效，在时钟沿采样输入；若写使能无效，保持原值。

\begin{lstlisting}
on clock edge:
  if write_enable:
    Q = D
  else:
    Q holds
\end{lstlisting}

寄存器的意义不只是“能存数”。它把连续传播的组合逻辑切成离散步骤。上一个周期的结果在寄存器输出端稳定存在，成为本周期组合逻辑的输入；本周期算出的结果，要到下一个时钟沿才成为新的状态。

例如一个 8-bit 寄存器保存 \code{00101100}。只要写使能为 0，输入端怎样变化都不会改变 Q。写使能为 1 且时钟沿到来时，D 端的 8 个 bit 被一起采样。这里“一起”很重要：寄存器保存的是同一个边界上的一组 bit，不是一位一位随意变化。

\section{二、计数器是带反馈的寄存器}

计数器的下一值来自当前值加一：

\begin{lstlisting}
next = current + 1
on clock edge:
  current = next
\end{lstlisting}

这个结构里有反馈，但反馈不是直接绕回去，而是经过加法器，并且被时钟边界切开。若 current 为 3，本周期组合逻辑算出 next 为 4；只有下一个时钟沿到来，寄存器才保存 4。随后组合逻辑再从 4 算出 5。

如果没有寄存器边界，current+1 的结果会立刻回到输入，再加一，再回到输入，电路无法表达“每个周期只增加一次”。计数器因此是理解同步电路的好例子：组合逻辑决定下一值，寄存器决定何时接受下一值。

计数器还会遇到宽度限制。4-bit 计数器从 \code{0000} 到 \code{1111}，再加一会回到 \code{0000}。这不是故障，而是硬件位宽决定的结果。后面讲溢出时会再次看到这个现象。

\section{三、移位寄存器移动 bit}

移位寄存器把一组 bit 向左或向右移动。它可以用于串行输入、串行输出、简单延迟线和位级变换。

\begin{lstlisting}
next[3] = current[2]
next[2] = current[1]
next[1] = current[0]
next[0] = serial_in
\end{lstlisting}

假设 current 为 \code{1010}，serial\_in 为 1。下一个时钟沿后，寄存器变成 \code{0101}。注意，四个位不是按顺序一个一个搬。组合逻辑同时准备好四个 next 输入，时钟沿到来时一起写入。

这个例子说明状态更新不一定是加法，也可以是重新排列、选择、清零、保持或装载外部输入。寄存器提供状态位置，组合逻辑定义下一状态规则。

\section{四、状态机把阶段保存下来}

有限状态机由三部分组成：当前状态、输入条件、下一状态逻辑。当前状态保存在寄存器里，组合逻辑根据当前状态和输入条件决定下一状态。

以一个三阶段硬件控制器为例：

\begin{longtable}{p{0.2\linewidth}p{0.32\linewidth}p{0.34\linewidth}}
\toprule
当前状态 & 条件 & 下一状态 \\
\midrule
IDLE & start=1 & LOAD \\
LOAD & ready=1 & RUN \\
RUN & count\_done=1 & DONE \\
DONE & clear=1 & IDLE \\
\bottomrule
\end{longtable}

若当前状态是 IDLE 且 start=0，状态保持 IDLE；若 start=1，下一状态逻辑准备 LOAD，等时钟沿到来后状态寄存器才真正变成 LOAD。状态机因此不会在一个周期里把 IDLE、LOAD、RUN、DONE 全部跑完，它每次只跨过一个明确边界。

\begin{pdfdiagram}[x=1cm,y=1cm]
  \node[hw state,minimum width=1.6cm] (idle) at (1.1,1.7) {IDLE};
  \node[hw state,minimum width=1.6cm] (load) at (4.0,1.7) {LOAD};
  \node[hw state,minimum width=1.6cm] (run) at (6.9,1.7) {RUN};
  \node[hw state,minimum width=1.6cm] (done) at (9.8,1.7) {DONE};
  \draw[hw bus] (idle.east) -- node[above,hw label] {start} (load.west);
  \draw[hw bus] (load.east) -- node[above,hw label] {ready} (run.west);
  \draw[hw bus] (run.east) -- node[above,hw label] {count\_done} (done.west);
  \draw[BookTeal,line width=0.55pt,-{Latex[length=2mm,width=1.5mm]}] (done.south) .. controls (7.8,0.2) and (3.2,0.2) .. node[below,hw label] {clear} (idle.south);
  \node[hw label,align=center] at (5.45,2.55) {当前状态保存在状态寄存器里\\下一状态由组合逻辑准备，时钟沿才提交};
\end{pdfdiagram}
\diagramnote{状态机每个周期只从当前状态跨到下一状态，不能在一个周期内把所有圆圈跑完。}

状态机输出可以只由当前状态决定，也可以同时依赖输入。前者输出更稳定，后者反应更快但更容易受输入毛刺影响。这里先记住一条原则：状态机不是画几个圆圈，它是“状态寄存器 + 下一状态组合逻辑 + 输出组合逻辑”的硬件结构。

\section{五、状态机为控制器铺路}

现在硬件已经能够保存一组 bit、按周期加一、按周期移动、按条件切换阶段。下一步要把这些能力用于更大的动作组织：哪些寄存器在这个周期写入，ALU 做哪种运算，哪一路数据被选择，下一阶段去哪里。

这就是控制器的基础。控制器本质上也是组合逻辑和状态寄存器的配合，只是它输出的不是一个普通数据值，而是一组用来打开、关闭、选择其他硬件路径的控制信号。进入控制器之前，还需要先讲数的表示和算术电路。

\begin{classicnote}
本章的标注对象是“状态转移表”。只画状态圆圈不够，必须写清当前状态、输入条件、下一状态和输出。

\begin{itemize}
\item 纸上实验：为一个三状态控制器写表：IDLE 等待 start，RUN 计数，DONE 等待 clear。每一行都写出本周期输出和下一状态。
\item 位级证据：寄存器是一组并排触发器；计数器是寄存器输出经过加法器反馈到输入；移位寄存器是固定连线选择了相邻 bit。
\item 控制证据：状态机输出可以只依赖状态，也可以依赖状态和输入；后者反应快，但更容易把输入毛刺传给控制线。
\item 检查题：若状态编码里出现未使用编码，电路上电或受干扰进入它时，下一状态逻辑应该怎样让机器回到安全状态？
\end{itemize}

经典系统教材会把状态机当作小型机器来读：状态寄存器是“当前 PC”，下一状态逻辑是“控制转移”，输出逻辑是“本周期动作”。理解这点，后面控制器和 CPU 就不再神秘。
\end{classicnote}
